\section{Research Questions}
\label{sec:research-questions}

We investigate whether validation workflows already present in
machine-learning research repositories are sensitive to controlled changes in
reproducibility-relevant experimental choices. The study addresses two
research questions.

\paragraph{RQ1.}
When reproducibility-relevant mutations are introduced into ML research
software, how often are they detected by existing repository validation
workflows?

\paragraph{RQ2.}
How does mutation detection differ by validation-workflow type and by the
strength of the workflow's validation oracle?

RQ1 concerns the overall sensitivity of the selected validation workflows to
confirmed non-equivalent mutations. RQ2 examines this sensitivity
descriptively across workflow and oracle categories that were fixed
independently of mutation outcomes.

Here, oracle strength is operationalized through the prospectively recorded
oracle categories described in Section~\ref{sec:oracle-classification}; these
categories are not assumed to form an ordinal scale.

Neither research question treats mutation survival as evidence that a
repository is irreproducible. A surviving mutation indicates only that the
selected validation workflow did not detect that particular controlled change.
